\subsection{Persistencia regional y unicidad del selector de acción}
\label{sec:accion-selector-regional}

La microfibra de clausura de la sección~\ref{sec:generacion} proporciona
el antecedente regional del lector de acción. Con la reducción ternaria orientada
\(\Psi:\mathcal U_6\to\mathbb F_3^6\) ya construida, sea
\(\mathcal F_\pi=\Psi^{-1}(010211)\). Sus cinco emisiones, separadas
en bloque inicial y terminal, son
\[
\begin{aligned}
 &(501,614,498\mid169,272,272),\\
 &(810,923,870\mid169,272,272),\\
 &(810,983,810\mid169,272,272),\\
 &(870,923,810\mid169,272,272),\\
 &(870,923,810\mid418,674,521).
\end{aligned}
\]
La proyección de retorno es
\[
 p_{\rm ret}:\mathbb Z^6\longrightarrow\mathbb Z^3,
 \qquad p_{\rm ret}(u_1,\ldots,u_6)=(u_4,u_5,u_6).
\]
En este paso cada emisión cuenta una vez: se selecciona la persistencia
entre las cinco regiones, conservando por separado las multiplicidades
de semillas que originan esas emisiones.

\begin{proposition}[Bloque terminal persistente]
\label{prop:accion-bloque-persistente}
El multiconjunto \(p_{\rm ret}(\mathcal F_\pi)\) tiene un único modo,
\[
 b_\pi=(169,272,272),\qquad \operatorname{mult}(b_\pi)=4.
\]
El bloque restante, \(b_\pi^-=(418,674,521)\), tiene multiplicidad uno.
\end{proposition}
\begin{proof}
La proyección de las primeras cuatro emisiones es \(b_\pi\); la de la
quinta es \(b_\pi^-\). Ambos bloques son distintos. La multiplicidad
máxima es, por tanto, cuatro y se alcanza en un solo bloque.
\end{proof}

El selector se define sobre multiconjuntos de bloques terminales con un
modo único de tipo \((r,s,s)\), \(r\ne s\). El grupo
\(\mathfrak S_3\) actúa renumerando las posiciones de esos bloques;
esta equivarianza compara sus representaciones, sin añadir una simetría
dinámica a la fibra fija \(\Psi^{-1}(010211)\).
El estabilizador de \(b_\pi\) intercambia las posiciones ocupadas por
\(272\) y fija la posición ocupada por \(169\). Para un bloque de tipo
\((r,s,s)\), con \(r\ne s\), el valor de esa posición fija es
\[
 \operatorname{sing}(b)=\sum_{j=1}^3 b_j-2\operatorname{moda}(b).
\]
La selección del modo regional seguida de esta operación define el
selector residual \(\rho_\pi\).

\begin{theorem}[Unicidad equivariante del retorno]
\label{thm:accion-selector-169}
Entre los selectores invariantes bajo permutaciones de las cinco regiones,
equivariantes bajo permutaciones de las posiciones terminales, que
seleccionan el bloque de multiplicidad máxima y retienen una posición
fija por su estabilizador, el valor de retorno es único:
\[
 \rho_\pi=\operatorname{sing}(b_\pi)=169.
\]
\end{theorem}
\begin{proof}
La invariancia respecto de las regiones hace depender la selección del
multiconjunto, no del orden de enumeración. La persistencia modal
selecciona \(b_\pi\) por la proposición~\ref{prop:accion-bloque-persistente}.
Su estabilizador tiene dos órbitas de posiciones: un par y un punto fijo.
La condición de posición fija selecciona este último, cuyo valor es
\(169\). Al permutar las posiciones, la equivarianza transporta esa
posición y conserva su valor. La construcción descrita satisface las
cuatro condiciones, lo que prueba también la existencia.
\end{proof}

La diferencia orientada entre ambos retornos conserva un dato adicional:
\[
 \omega_\pi=b_\pi^- - b_\pi=(249,402,249),\qquad
 \langle(1,1,1),\omega_\pi\rangle=900.
\]
El selector \(169\) es un dato terminal de grado cero. La diferencia
\(\omega_\pi\) asigna un dato de grado uno a la arista orientada entre
los dos bloques; su tratamiento como cociclo exige especificar el complejo
y comprobar el coborde correspondiente. Esta distinción conserva tanto
el retorno como su variación, sin utilizarlos indistintamente.

La selección es regional: en el catálogo completo de \(468\) emisiones,
\(169\) aparece \(84\) veces, \(14\) en cada posición. Su función en
el lector procede de la microfibra, de la persistencia modal y del
estabilizador, conjuntamente; la frecuencia global se conserva como otro
recuento del mismo catálogo.

\subsection{Longitud de palabra y carácter de retorno definido por un determinante}
\label{sec:accion-grado-determinante}

Sea \(\mathscr P_{\rm pal}=\bigoplus_{n\ge0}\mathscr P_{{\rm pal},n}\) el módulo libre
sobre las palabras finitas del alfabeto de emisión, graduado por longitud,
con la filtración inducida. Este dominio contiene las emisiones regionales
y sus concatenaciones; la concatenación suma grados. El carácter
formal de longitud
\[
 \chi_z([w])=z^{|w|},\qquad
 \chi_z([uv])=\chi_z([u])\chi_z([v])
\]
conserva esta operación. La compatibilidad cilíndrica del lector envía
\(\mathscr P_{{\rm pal},n}\) al componente homogéneo \(z^n\mathbb R\) de
\(\mathbb R[z]\); la evaluación \(z=\alpha\) se efectúa después.
Como cada \(w\in\mathcal F_\pi\) tiene longitud seis, el impulso
residual es independiente de la región elegida y tiene grado seis:
\[
 \rho_\pi\chi_z([w])=169z^6,
 \qquad\left.\rho_\pi\chi_z([w])\right|_{z=\alpha}=169\alpha^6.
\]
La selección registra un impulso común; no suma de nuevo las cinco
emisiones después de haber extraído su retorno persistente.
La longitud de la emisión y la longitud de un paseo en el grafo de
incidencias son grados de objetos diferentes. En particular, los paseos
cerrados que recorren los cuatro anclajes regionales tienen longitud
mínima ocho; el grado utilizado aquí cuenta las seis posiciones de una
emisión, no las aristas de ese paseo.

El cierre dodecafásico y la completación de las tres parejas nonádicas
proporcionan la coordenada común \(A=10^3\alpha\). En radianes,
\[
 x=\frac{\pi A}{180}=\frac{50\pi}{9}\alpha.
\]
Sea \(R_{\rm h}\) una involución real bidimensional de traza nula,
\(R_{\rm h}^2=I_2\), \(\operatorname{tr}R_{\rm h}=0\). El transporte bicapa es
\[
 \mathcal T(x,y)=\exp(-xI_2-yR_{\rm h}).
\]
La coordenada orientada \(y\) queda sin evaluar. Como el polinomio
\(X^2-1\) tiene raíces distintas y la traza de \(R_{\rm h}\) es cero,
sus autovalores son \(1,-1\). En una base de autovectores, los dos
factores del transporte son \(e^{-x-y}\) y \(e^{-x+y}\). Por tanto,
\[
 \det\mathcal T=e^{-x-y}e^{-x+y}=e^{-2x}
 =\exp\!\left(-\frac{100\pi\alpha}{9}\right)=D_A.
\]
Sobre la línea determinante, \(\bigwedge^2\mathcal T\) actúa como
\(D_A\) por la identidad. El producto elimina la coordenada orientada y conserva la contracción
común. El determinante es invariante bajo conjugación y bajo el
intercambio de hojas \(R_{\rm h}\mapsto-R_{\rm h}\); así se construye \(D_A\)
antes de evaluar el contraángulo.

\begin{proposition}[Carácter multiplicativo del retorno]
\label{prop:accion-caracter-determinantal}
El carácter multiplicativo \(\delta_A:(\mathbb N_0,+)\to(\mathbb R_{>0},\cdot)\)
cuyo valor en una prolongación elemental es \(D_A\) queda determinado
por \(\delta_A(k)=D_A^k\).
\end{proposition}
\begin{proof}
Un carácter preserva la unidad, por lo que \(\delta_A(0)=1\).
La concatenación de \(k\) prolongaciones da
\(\delta_A(k)=\delta_A(1)^k=D_A^k\). Estas potencias satisfacen
la multiplicatividad y el valor elemental prescrito, probando existencia
y unicidad.
\end{proof}

\Needspace{17\baselineskip}
\begin{definition}[Reglas de la continuación regional]
\label{def:accion-lector-regional}
El lector regional definido mediante el determinante del transporte conserva cinco operaciones:
\begin{enumerate}
 \item la longitud cilíndrica determina el grado analítico;
 \item el impulso de retorno y la continuación estricta son sumandos
 distintos de una renovación;
 \item tras registrar \(\rho_\pi\) una sola vez, la continuación
 comienza en la unidad de la línea determinante;
 \item cada prolongación actúa por \(\bigwedge^2\mathcal T\), cuyo
 factor escalar es \(D_A\);
 \item en la orientación cardinal de APP, el retorno regional ocupa el
 sector compensador y se sustrae de la fase de acción de quinto grado.
\end{enumerate}
\end{definition}

La normalización unitaria determina la amplitud de la continuación, además
de su razón. En efecto, todas las series
\[
 F_\lambda(z)=169z^6+\lambda\frac{D_Az^7}{1-D_Az}
 \qquad(\lambda\in\mathbb R)
\]
conservan el impulso finito y la misma recurrencia multiplicativa de los
coeficientes de la cola. La tercera regla fija \(\lambda=1\) antes de evaluar
\(z=\alpha\). El residuo se registra una vez; la prolongación cuenta
la línea que continúa, con su unidad. La ecuación de renovación
\eqref{eq:revision-renovacion} y su prueba de unicidad, desarrolladas
a continuación, completan esta cadena entre región, grado,
transporte y sección de acción.
